%% kapa: exam=1 round=26 year=2015 period=2 form=A subject=appraisal_law subject_ko=부동산관계법규 seq=23 number=23
%% source: Q-Net 감정평가사 기출문제 (cst00309) artlSeq=5138001 file=extracted/1008/5138001/제26회_감정평가사 1차_기출/제26회_감정평가사 1차_기출/2교시_감평_A형.pdf
%% conversion: kordoc 4.24.51 -> tools/gen1.py (auto); review=proofread-sonnet; flags=none
\begin{question}{23}
부동산가격공시 및 감정평가에 관한 법령상 표준지공시지가에 관한 설명으로 옳은 것은?
\begin{choices}
  \item 지역분석조서는 표준지조사평가보고서에 첨부되지 않는다.
  \item 국토교통부장관은 표준지공시지가에 대한 이의신청의 내용이 타당하다고 인정될 때에는 당해 표준지공시지가를 조정하되, 이를 다시 공시하여야 하는 것은 아니다.
  \item 국토교통부장관은 토지가격 산정의 목적을 위한 지가산정을 위하여 필요하다고 인정하는 경우에는 토지가격비준표를 작성하여 관계행정기관 등에 제공하여야 한다.
  \item 표준지공시지가에 대한 이의신청에 관한 사항을 심의하기 위하여 국토교통부장관 소속으로 중앙토지수용위원회를 둔다.
  \item 표준지공시지가의 공시사항에는 표준지 및 주변토지의 이용상황은 포함되지 않는다.
\end{choices}
\end{question}
\answer{3}
